\documentclass[
    %aspectratio=169,
]{beamer}

\usepackage[english]{babel}
\usepackage[utf8]{inputenc}
\usepackage[T1]{fontenc}
\usepackage{csquotes}
\usepackage{expl3,biblatex}

\addbibresource{bibliography.bib}

\usepackage{booktabs}
\usetheme{Madrid}
\usecolortheme{crane}
\usepackage{svg}
\usepackage{graphicx,caption}

% set font
\usefonttheme{serif}

% to prevent the boxes from overlapping the logo at the lower right corner
\addtobeamertemplate{block begin}{%
  \setlength{\textwidth}{0.9\textwidth}%
}{}

\addtobeamertemplate{block alerted begin}{%
  \setlength{\textwidth}{0.9\textwidth}%
}{}

\addtobeamertemplate{block example begin}{%
  \setlength{\textwidth}{0.9\textwidth}%
}{}


% usage: \hl{text to emphasis as yellow background}
\usepackage{soul}
\makeatletter
\let\HL\hl
\renewcommand\hl{%
  \let\set@color\beamerorig@set@color
  \let\reset@color\beamerorig@reset@color
  \HL}
\makeatother


% usage: create outline pages before each section
\AtBeginSection[]
{
  \begin{frame}
    \frametitle{Table of Contents}
    \tableofcontents[currentsection]
  \end{frame}
}

% Title page information
\title[Short Presentation Title]{Full Presentation Title}
\subtitle[Short Subtitle]{Full Subtitle}
\author[Short Name or Initials]{Full Author Name\texorpdfstring{\\}{, }Author1234@vandals.uidaho.edu}
\institute[U of I]{University of Idaho}
\date{\today}
\logo{\includesvg[height=1cm]{./uibeamer/logo/ui-stacked-gold-black.svg}}
\subject{Presentation Subject}
\keywords{the, presentation, keywords}

\begin{document}

% make the title page
\begin{frame}%[plain]
\maketitle
\end{frame}

% make the outline page/table of contents
\begin{frame}
\frametitle{Table of Contents}
\tableofcontents
\end{frame}
\section[Short Section 1 Name]{Full Section 1 Name}
\subsection[Short Subsection 1 Name]{Full Subsection 1 Name}

\begin{frame}{Frame Title}{Frame Subtitle}
plain text, \structure{page structure}, \alert{emphasis}, \hl{high light emphasis}
\begin{itemize}
  \item a single-line bullet list item
  \item a bullet list item that is quite long (in order to force a line break),
    which also contains \alert{emphasized text}
  \begin{itemize}
    \item a second-level list item
    \begin{itemize}
      \item a third-level list item
    \end{itemize}
    \item \alert{an emphasized second-level list item}
  \end{itemize}
\end{itemize}
\begin{enumerate}
  \item a numbered list item
  \begin{enumerate}
    \item a second-level list item containing a math expression
      \[ E = mc^2 \]
      and a citation \cite{einstein1905tragheit}
  \end{enumerate}
\end{enumerate}
\end{frame}


\subsection[Short Subsection 2 Name]{Full Subsection 2 Name}

\begin{frame}{Text Blocks}
text above a block\footnote{a footnote with an \url{https://address.edu}}
\begin{block}{Block}
  text in a block
\end{block}
\begin{exampleblock}{Example Block}
  text in a block
\end{exampleblock}
\begin{alertblock}{Emphasized Block}
  text in a block
\end{alertblock}
\end{frame}


\begin{frame}{Figures}{"Potatoes."}

\begin{figure}
    \includegraphics[width=.5\textwidth,height=.5\textheight,keepaspectratio]{./figures/a-sketch-of-a-dish-of-potato-from-Idaho-state.jpg}
      \caption{a sketch of a dish of potato from Idaho state%
        \footnote{%
          Do you like famous potatoes?
          They do.
          \cite{evelyn1984potatoes}
        }%
      }

\end{figure}
\end{frame}


\subsection[Short Subsection 3 Name]{Full Subsection 3 Name}


\begin{frame}{Sample Data}
\begin{table}
  \begin{tabular}{llc}
    Category & Item & Quantity \\ \midrule
    A & Item A1 & 15 \\
    A & Item A2 & 20 \\
    B & Item B1 & 25 \\
    B & Item B2 & 30 \\
    C & Item C1 & 10 \\
  \end{tabular}
  \caption{Example data for demonstration purposes}
\end{table}
\end{frame}


\section[Equation]{Equation and Explanations}

\begin{frame}{Equation and Explanations}{sample equation}
\textbf{Functionality Representation Model (FRM)} is an abstract method that aims to analyze the significance of various components in a complex system.

$$f(z) = \alpha + \sum^{N}_{i=1}\beta_iz_i$$

In this equation, $z$ symbolizes a data point, $z_i$ denotes the $i^{th}$ component of $z$, and $\beta_i$ represents the influence of component $z_i$ on the system's behavior.
    
\end{frame}

\section{\bibname}
\begin{frame}[t, allowframebreaks]{\bibname}
\printbibliography[heading=none]
\end{frame}


%\miniframesoff


\begin{frame}%[plain]
\vfill
\centerline{Thank You for Your Attention!}
\vfill\vfill
\end{frame}
%\input{sample}

\end{document}


